\documentclass[activite]{thedocument}
\startdatakeys
\source{}
\cycle{}
\competencesDuSocle{}
\connaissancesRequises{}
\competencesTravaillees{}
\enddatakeys
\begin{document}
    \begin{enumerate}
        \item Calculer l'aire du rectangle ABCD de deux manières différentes.
        \begingroup\centering
        \begin{tikzpicture}[scale=2]
            \coordinate[label=below left:A](A)at(0,0);
            \coordinate[label=below right:B](B)at(3,0);
            \coordinate[label=above right:C](C)at(3,2);
            \coordinate[label=above left:D](D)at(0,2);
            \coordinate(E1)at(2,0);
            \coordinate(E2)at(2,2);
            \coordinate(F1)at(0,0.5);
            \coordinate(F2)at(3,0.5);
            \draw(A)--(B)--(C)--(D)--cycle;
            \draw(F1)--(F2);
            \draw(A)--(B)node[pos=0.5,below]{6cm};
            \draw(B)--(F2)node[pos=0.5,right]{1 cm};
            \draw(F2)--(C)node[pos=0.5,right]{2 cm};
        \end{tikzpicture}\par\vskip50pt\endgroup
        \item%
        Écrire deux expressions littérales donnant l'aire du rectangle ABCD
        ci-dessous.\vskip30pt
        \begingroup\centering
        \begin{tikzpicture}[scale=2]
            \coordinate[label=below left:A](A)at(0,0);
            \coordinate[label=below right:B](B)at(3,0);
            \coordinate[label=above right:C](C)at(3,2);
            \coordinate[label=above left:D](D)at(0,2);
            \coordinate(E1)at(2,0);
            \coordinate(E2)at(2,2);
            \coordinate(F1)at(0,0.5);
            \coordinate(F2)at(3,0.5);
            \draw(A)--(B)--(C)--(D)--cycle;
            \draw(F1)--(F2);
            \draw(A)--(B)node[pos=0.5,below]{$k$};
            \draw(B)--(F2)node[pos=0.5,right]{$a$};
            \draw(F2)--(C)node[pos=0.5,right]{$b$};
        \end{tikzpicture}\par\vskip50pt\endgroup
        \item En déduire une formule liant les nombres $k$, $a$ et $b$.
    \end{enumerate}
\end{document}